\chapter{Vektor, Garis dan Bidang di Ruang}\label{ch:g12:space}

Geometri tiga dimensi menjadi dapat dihitung begitu \omterm{def:g11:vect:vector}{vektor} dan
\omterm{def:g10:coordgeom:system}{koordinatnya} tersedia: garis memperoleh perwakilan parametrik, bidang memperoleh
\omterm{def:g10:algebra:equation}{persamaan} Kartesius, dan \omterm{def:g12:space:dot}{hasil kali skalar} mengukur sudut dan jarak.
Bab ini membangun kotak perkakas itu lalu memakainya untuk menyelesaikan soal
perpotongan dan jarak yang klasik.

\section{Vektor di ruang}

\omterm{def:g11:vect:vector}{Vektor} di ruang menuruti kaidah yang sama seperti di bidang: keduanya menjumlah,
menskala, dan $\vect{AB} = \vect{CD}$ tepat ketika $ABDC$ jajaran genjang. Pada suatu
\omterm{def:g10:coordgeom:system}{sistem koordinat} $(O; \vec\imath, \vec\jmath, \vec k)$, setiap \omterm{def:g11:vect:vector}{vektor} mempunyai
\omterm{def:g10:coordgeom:system}{koordinat} $(x, y, z)$.

\begin{definition}[Kesegarisan, kesebidangan]\label{def:g12:space:collinear}
Dua \omterm{def:g11:vect:vector}{vektor} disebut \emph{segaris} jika yang satu kelipatan yang lain. Tiga
\omterm{def:g11:vect:vector}{vektor} $\vec u, \vec v, \vec w$ disebut \emph{sebidang}\index{sebidang} jika salah
satunya dapat ditulis sebagai gabungan dua \omterm{def:g11:vect:vector}{vektor} lainnya, misalnya
$\vec w = a\vec u + b\vec v$.
\end{definition}

\begin{proposition}\label{prop:g12:space:basis}
Jika $\vec u, \vec v, \vec w$ tidak \omterm{def:g12:space:collinear}{sebidang}, maka setiap \omterm{def:g11:vect:vector}{vektor} di ruang
terurai secara tunggal sebagai $x\vec u + y\vec v + z\vec w$.
\end{proposition}

\begin{proof}[Gagasan buktinya]
Lewat ujung \omterm{def:g11:vect:vector}{vektor} yang hendak diuraikan, lukislah garis yang sejajar dengan
$\vec w$; garis itu menemui bidang $\vec u, \vec v$ pada satu titik saja, dan titik
itu memecah \omterm{def:g11:vect:vector}{vektornya} menjadi komponen di bidang tersebut (secara tunggal
$x\vec u + y\vec v$, menurut geometri bidang) dan komponen sepanjang $\vec w$.
Ketunggalannya: selisih dua penguraian akan menyatakan salah satu \omterm{def:g11:vect:vector}{vektornya}
dengan dua \omterm{def:g11:vect:vector}{vektor} lainnya, dan itu bertentangan dengan ketaksebidangannya.
\end{proof}

\section{Garis dan bidang}

\begin{definition}[Perwakilan parametrik sebuah garis]\label{def:g12:space:line}
Garis lewat $A(x_A, y_A, z_A)$ dengan \omterm{def:g11:vect:direction}{vektor arah}
$\vec u (a, b, c) \neq \vec 0$ adalah himpunan titik
\[
M(x_A + ta,\; y_A + tb,\; z_A + tc), \qquad t \in \R .
\]
\end{definition}

\begin{definition}[Bidang]\label{def:g12:space:plane}
Bidang lewat $A$ yang diarahkan oleh dua \omterm{def:g11:vect:vector}{vektor} yang tidak \omterm{def:g12:space:collinear}{segaris},
$\vec u, \vec v$, adalah himpunan titik $M$ dengan
$\vect{AM} = s\vec u + t\vec v$, $(s, t) \in \R^2$.
\end{definition}

\begin{omfigure}
\begin{tikzpicture}[scale=1.05]
  \draw (-0.5,-0.25) -- (3.4,1.7);
  \draw[-Stealth, omDef, thick] (1,0.5) -- (1.9,0.95)
    node[midway, above left, font=\small] {$\vec u$};
  \fill (0.1,0.05) circle (1.2pt) node[below right, font=\small] {$t=-1$};
  \fill (1,0.5) circle (1.2pt) node[above left, font=\small] {$A$};
  \fill (1.9,0.95) circle (1.2pt) node[below right, font=\small] {$t=1$};
  \fill (2.8,1.4) circle (1.2pt) node[below right, font=\small] {$t=2$};
\end{tikzpicture}
\qquad\quad
\begin{tikzpicture}[scale=1.05]
  \draw (0,0) -- (4,0.5) -- (5.2,1.9) -- (1.2,1.4) -- cycle;
  \node[font=\small] at (4.75,1.6) {$\mathcal P$};
  \coordinate (A) at (1.3,0.55);
  \coordinate (M) at (3.77,1.38);
  \draw[-Stealth, omDef, thick] (A) -- ++(1.31,0.16)
    node[midway, below, font=\small] {$\vec u$};
  \draw[-Stealth, omDef, thick] (A) -- ++(0.42,0.49)
    node[midway, left, font=\small] {$\vec v$};
  \draw[black, dashed, thin] (3.27,0.79) -- (M) -- (1.80,1.14);
  \fill (A) circle (1.2pt) node[below left, font=\small] {$A$};
  \fill (M) circle (1.2pt) node[above, font=\small] {$M$};
\end{tikzpicture}

{\small Kiri: garis lewat $A$ dengan arah $\vec u$, yang berskala menurut
parameter $t$. Kanan: bidang lewat $A$ yang diarahkan oleh $\vec u$ dan
$\vec v$; setiap titik $M$ pada bidang itu dicapai sebagai
$\vect{AM} = s\vec u + t\vec v$.}
\end{omfigure}

\begin{proposition}[Kedudukan relatif]\label{prop:g12:space:positions}
Dua bidang yang berbeda entah sejajar entah berpotongan pada sebuah garis. Sebuah
garis dan sebuah bidang entah sejajar (mungkin termuat) entah berpotongan pada satu
titik. Dua garis di ruang dapat berpotongan, sejajar secara tegas, berimpit,
atau \emph{bersilangan} (tidak \omterm{def:g12:space:collinear}{sebidang}).
\end{proposition}

\begin{proof}[Sketsa]
Semuanya pernyataan tentang \omterm{def:g11:prob:model}{kejadian}; setiap penelaahan kasusnya menyusut, dalam
\omterm{def:g10:coordgeom:system}{koordinat}, menjadi penyelesaian sistem linear lalu pencacahan penyelesaiannya --- lihat
\cref{met:g12:space:intersections}.
\end{proof}

\section{Hasil kali skalar}

\begin{definition}[Hasil kali skalar]\label{def:g12:space:dot}
\emph{Hasil kali skalar}\index{hasil kali skalar} dari $\vec u$ dan $\vec v$ adalah
\[
\vec u \cdot \vec v = \tfrac12\left(\norm{\vec u + \vec v}^2
- \norm{\vec u}^2 - \norm{\vec v}^2\right),
\]
dengan $\norm{\vec u}$ panjang $\vec u$. Jika kedua \omterm{def:g11:vect:vector}{vektornya}
tak nol, maka $\vec u \cdot \vec v = \norm{\vec u}\,\norm{\vec v}\cos\theta$
dengan $\theta$ sudut di antara keduanya; dan pada \omterm{def:g10:coordgeom:system}{sistem koordinat} yang
\emph{\omterm{def:g10:coordgeom:system}{ortonormal}},
\[
\vec u \cdot \vec v = xx' + yy' + zz' .
\]
Dua \omterm{def:g11:vect:vector}{vektor} disebut \emph{\omterm{def:g11:scal:orthogonal}{ortogonal}} jika hasil kali skalarnya $0$.
\end{definition}

\begin{proposition}[Sifatnya]\label{prop:g12:space:dotprops}
\omterm{def:g12:space:dot}{Hasil kali skalar} bersifat setangkup ($\vec u\cdot\vec v = \vec v\cdot\vec u$),
bilinear ($(a\vec u + b\vec v)\cdot \vec w = a\,\vec u\cdot\vec w +
b\,\vec v\cdot \vec w$), dan
$\vec u \cdot \vec u = \norm{\vec u}^2 \geq 0$.
\end{proposition}

\begin{proof}
Pada \omterm{def:g10:coordgeom:system}{sistem ortonormal}, ketiga sifatnya langsung terlihat pada rumus
$xx' + yy' + zz'$, yang sendirinya menyusul dari rumus pendefinisinya dan
perhitungan panjang lewat Pythagoras:
$\norm{\vec u}^2 = x^2 + y^2 + z^2$.
\end{proof}

\begin{definition}[Vektor normal]\label{def:g12:space:normal}
Suatu \emph{vektor normal}\index{vektor normal} bagi bidang $\mathcal P$ adalah
\omterm{def:g11:vect:vector}{vektor} tak nol yang \omterm{def:g11:scal:orthogonal}{ortogonal} terhadap setiap \omterm{def:g11:vect:vector}{vektor} yang terletak pada $\mathcal P$
(dan cukuplah ia \omterm{def:g11:scal:orthogonal}{ortogonal} terhadap dua arah $\mathcal P$ yang tidak
\omterm{def:g12:space:collinear}{segaris}).
\end{definition}

\begin{omfigure}
\begin{tikzpicture}[scale=1.1]
  \draw (0,0) -- (4,0.5) -- (5.2,1.9) -- (1.2,1.4) -- cycle;
  \node[font=\small] at (4.75,1.6) {$\mathcal P$};
  \coordinate (H) at (2.6,0.95);
  \coordinate (M0) at (2.86,3.03);
  \draw[-Stealth, omThm, thick] (H) -- (2.73,1.99)
    node[midway, right=1pt, font=\small] {$\vec n$};
  \draw[black, dashed, thin] (2.73,1.99) -- (M0);
  \draw[thin] (2.79,0.974) -- (2.815,1.174) -- (2.625,1.15);
  \draw[black, dashed, thin] (H) -- (1.1,0.55);
  \fill (H) circle (1.2pt) node[below right, font=\small] {$H$};
  \fill (M0) circle (1.2pt) node[above, font=\small] {$M_0$};
  \fill (1.1,0.55) circle (1.2pt) node[below left, font=\small] {$A$};
\end{tikzpicture}

{\small \omterm{def:g12:space:normal}{Vektor normal} $\vec n$ bagi bidang $\mathcal P$ \omterm{def:g11:scal:orthogonal}{ortogonal}
terhadap setiap arah $\mathcal P$. Titik $H$, yaitu kaki
tegak lurus dari $M_0$, mewujudkan jarak $M_0$ ke bidangnya.}
\end{omfigure}

\begin{theorem}[Persamaan Kartesius sebuah bidang]\label{thm:g12:space:planeeq}
Pada \omterm{def:g10:coordgeom:system}{sistem koordinat} \omterm{def:g10:coordgeom:system}{ortonormal}, bidang lewat
$A$ dengan \omterm{def:g12:space:normal}{vektor normal} $\vec n(a, b, c)$ berpersamaan
\[
a x + b y + c z + d = 0,
\]
dengan $d = -(ax_A + by_A + cz_A)$. Sebaliknya, setiap \omterm{def:g10:algebra:equation}{persamaan} semacam itu dengan
$(a,b,c) \neq (0,0,0)$ mendefinisikan bidang dengan \omterm{def:g12:space:normal}{vektor normal} $(a, b, c)$.
\end{theorem}

\begin{proof}
$M \in \mathcal P \iff \vect{AM} \perp \vec n \iff
a(x - x_A) + b(y - y_A) + c(z - z_A) = 0$, dan itu terjabar menjadi \omterm{def:g10:algebra:equation}{persamaannya}.
Sebaliknya, ambillah sebarang titik penyelesaian $A$ dari \omterm{def:g10:algebra:equation}{persamaannya}; perhitungan
yang sama secara mundur menunjukkan bahwa himpunan penyelesaiannya adalah
$\{M : \vect{AM}\cdot\vec n = 0\}$, yaitu bidang lewat $A$ yang normalnya
$\vec n$.
\end{proof}

\begin{method}[Perpotongan dalam praktiknya]\label{met:g12:space:intersections}
\begin{itemize}
  \item \emph{Garis $\cap$ bidang}: substitusikan \omterm{def:g10:coordgeom:system}{koordinat} parametrik
        garisnya ke dalam \omterm{def:g10:algebra:equation}{persamaan} bidangnya; lalu selesaikan untuk $t$ (satu
        penyelesaian: satu titik; $0 = \text{tak nol}$: sejajar; $0 = 0$: garisnya termuat).
  \item \emph{Bidang $\cap$ bidang}: selesaikan sistem kedua \omterm{def:g10:algebra:equation}{persamaannya},
        dengan memarameterkan lewat satu \omterm{def:g10:coordgeom:system}{koordinat} bebas; penyelesaiannya sebuah garis
        (atau bidangnya sejajar ketika \omterm{def:g12:space:normal}{vektor normalnya} \omterm{def:g12:space:collinear}{segaris}).
  \item \emph{Pemeriksaan keortogonalan}: garis $\perp$ bidang jika dan hanya jika
        arahnya \omterm{def:g12:space:collinear}{segaris} dengan normalnya; dan dua bidang tegak lurus jika dan hanya
        jika normalnya \omterm{def:g11:scal:orthogonal}{ortogonal}.
\end{itemize}
\end{method}

\begin{proposition}[Jarak titik ke bidang]\label{prop:g12:space:distance}
\index{jarak!titik ke bidang}
Pada \omterm{def:g10:coordgeom:system}{sistem ortonormal}, jarak dari $M_0(x_0, y_0, z_0)$ ke
bidang $\mathcal P : ax + by + cz + d = 0$ adalah
\[
\operatorname{dist}(M_0, \mathcal P)
= \frac{\abs{ax_0 + by_0 + cz_0 + d}}{\sqrt{a^2 + b^2 + c^2}}.
\]
\end{proposition}

\begin{proof}
Misalkan $H$ proyeksi \omterm{def:g11:scal:orthogonal}{ortogonal} $M_0$ pada $\mathcal P$, yaitu kaki
tegak lurusnya, sehingga $\vect{HM_0}$ \omterm{def:g12:space:collinear}{segaris} dengan normal satuannya
$\frac{\vec n}{\norm{\vec n}}$, dan jaraknya adalah
$\abs{\vect{HM_0}\cdot \frac{\vec n}{\norm{\vec n}}}$. Untuk sebarang
$H \in \mathcal P$, berlaku
$\vect{HM_0}\cdot\vec n = ax_0 + by_0 + cz_0 - (ax_H + by_H + cz_H)
= ax_0 + by_0 + cz_0 + d$ (dengan memakai \omterm{def:g10:algebra:equation}{persamaannya} di $H$), dan dari situ
rumusnya sesudah dibagi $\norm{\vec n} = \sqrt{a^2+b^2+c^2}$.
\end{proof}

\begin{example}\label{ex:g12:space:distance}
Jarak dari titik asal ke bidang $x + 2y + 2z - 6 = 0$:
$\frac{\abs{-6}}{\sqrt{1 + 4 + 4}} = \frac63 = 2$.
\end{example}

\section{Latihan}

Pada semua latihannya \omterm{def:g10:coordgeom:system}{sistem koordinatnya} \omterm{def:g10:coordgeom:system}{ortonormal}.

\begin{exercise}[$\star$]\label{exo:g12:space:1}
Misalkan $A(1, 0, 2)$, $B(3, 1, 1)$ dan $C(2, -1, 3)$. Apakah titik $A$, $B$,
$C$ \omterm{def:g12:space:collinear}{segaris}? Berikan perwakilan parametrik bagi garis $(AB)$.
\end{exercise}

\begin{exercise}[$\star$]\label{exo:g12:space:2}
Tentukan \omterm{def:g10:algebra:equation}{persamaan} Kartesius bidang lewat $A(1, 1, 0)$ dengan
\omterm{def:g12:space:normal}{vektor normal} $\vec n(2, -1, 3)$. Apakah titik $B(0, 2, 1)$ terletak padanya?
\end{exercise}

\begin{exercise}[$\star$]\label{exo:g12:space:3}
Hitunglah sudut di $A$ (sampai derajat terdekat) pada segitiga yang bertitik
sudut $A(0,0,0)$, $B(1, 1, 0)$ dan $C(1, 0, 1)$.
\end{exercise}

\begin{exercise}[$\star\star$]\label{exo:g12:space:4}
Tinjaulah garis $\mathcal D$ lewat $A(1, 2, 0)$ dengan arah
$\vec u(1, -1, 2)$, dan bidang $\mathcal P : 2x + y - z + 1 = 0$.
Tentukan $\mathcal D \cap \mathcal P$.
\end{exercise}

\begin{exercise}[$\star\star$]\label{exo:g12:space:5}
Misalkan $\mathcal P : x - y + 2z = 1$ dan $\mathcal Q : 2x + y + z = 4$.
Tunjukkan bahwa $\mathcal P$ dan $\mathcal Q$ berpotongan pada sebuah garis lalu
berikan perwakilan parametriknya.
\end{exercise}

\begin{exercise}[$\star\star$]\label{exo:g12:space:6}
Tunjukkan bahwa garis
\[
\mathcal D_1 : (x, y, z) = (1 + t,\; 2t,\; 3 - t), \qquad
\mathcal D_2 : (x, y, z) = (2 + s,\; 1 + s,\; 1 + 2s)
\]
bersilangan (tidak \omterm{def:g12:space:collinear}{sebidang}).
\end{exercise}

\begin{exercise}[$\star\star$]\label{exo:g12:space:7}
Misalkan $\mathcal P : 2x - y + 2z - 3 = 0$.
\begin{enumerate}
  \item Hitunglah jarak dari $M_0(3, 1, 1)$ ke $\mathcal P$.
  \item Tentukan \omterm{def:g10:coordgeom:system}{koordinat} proyeksi \omterm{def:g11:scal:orthogonal}{ortogonal} $H$ dari
        $M_0$ pada $\mathcal P$, lalu periksalah jarak $M_0H$.
\end{enumerate}
\end{exercise}

\begin{exercise}[$\star\star\star$]\label{exo:g12:space:8}
Kubus $ABCDEFGH$ berusuk $1$; tempatkanlah \omterm{def:g10:coordgeom:system}{koordinatnya} sedemikian sehingga
$A(0,0,0)$, $B(1,0,0)$, $D(0,1,0)$, $E(0,0,1)$ (dengan
$C = B + D$, $F = B + E$, $G = B + D + E$, $H = D + E$ sebagai \omterm{def:g11:vect:vector}{vektor} dari
$A$).
\begin{enumerate}
  \item Tunjukkan bahwa diagonal $(AG)$ \omterm{def:g11:scal:orthogonal}{ortogonal} terhadap bidang $(BDE)$.
  \item Hitunglah jarak dari $A$ ke bidang $(BDE)$, dan titik
        tempat $(AG)$ menembusnya.
\end{enumerate}
\end{exercise}

\begin{exercise}[$\star\star\star$]\label{exo:g12:space:9}
\emph{(Volume bidang empat.)} Misalkan $A(1,1,1)$, $B(2,1,0)$, $C(0,2,1)$
dan $D(2,2,2)$.
\begin{enumerate}
  \item Periksalah bahwa $\vect{AB}$, $\vect{AC}$, $\vect{AD}$ tidak \omterm{def:g12:space:collinear}{sebidang}
        (sehingga titiknya membentuk bidang empat yang sungguhan).
  \item Carilah \omterm{def:g11:vect:vector}{vektor} $\vec n$ yang \omterm{def:g11:scal:orthogonal}{ortogonal} terhadap $\vect{BC}$ sekaligus
        $\vect{BD}$, lalu simpulkan \omterm{def:g10:algebra:equation}{persamaan} Kartesius bidang $(BCD)$.
  \item Hitunglah jarak dari $A$ ke bidang $(BCD)$ dan luas
        segitiga $BCD$, dengan memakai rumus
        $\text{Luas} = \frac12\sqrt{\norm{\vec u}^2\norm{\vec v}^2 -
        (\vec u \cdot \vec v)^2}$ bagi segitiga yang direntang
        $\vec u, \vec v$.
  \item Simpulkan volume bidang empatnya
        ($V = \frac13 \times \text{luas alas} \times \text{tinggi}$).
\end{enumerate}
\end{exercise}

\section{Soal: Mengiris kubus}

\begin{problem}[{Soal akhir pekan --- sebuah bidang memotong kubus menjadi
segi enam yang sempurna, sebuah bidang empat beraturan bersembunyi di sudutnya, dan
sudut diagonalnya membentuk setiap intan}]
\label{pb:g12:space:1}
Irislah sebuah kubus dan kamu akan menduga persegi dan persegi panjang --- padahal satu
\omterm{def:g10:numbers:interunion}{irisan} yang masyhur menghasilkan \emph{segi enam beraturan}, dan empat \omterm{def:g10:numbers:interunion}{irisan}
sudut yang cerdik meninggalkan \emph{bidang empat beraturan}. Kedua
kejutan itu, beserta sudut $109.5^\circ$ yang dipuja atom karbon,
tunduk pada kotak perkakas bab ini: garis parametrik, \omterm{def:g10:algebra:equation}{persamaan} bidang
dan \omterm{def:g12:space:dot}{hasil kali skalar} (\cref{thm:g12:space:planeeq},
\cref{prop:g12:space:distance}). Pakailah di seluruhnya kubus pada
\cref{exo:g12:space:8}: berusuk $1$, $A(0,0,0)$, $B(1,0,0)$,
$D(0,1,0)$, $E(0,0,1)$, $C(1,1,0)$, $F(1,0,1)$, $H(0,1,1)$,
$G(1,1,1)$.

\medskip
\noindent\textbf{Bagian I --- Kelancaran.}
\begin{enumerate}
  \item Tulislah \omterm{def:g10:algebra:equation}{persamaan} bidang lewat $B(1,0,0)$ dengan
        \omterm{def:g12:space:normal}{vektor normal} $(1,1,1)$, perwakilan parametrik bagi
        garis lewat $A$ yang diarahkan oleh $(1,1,1)$, lalu
        titik potongnya.
  \item Hitunglah jarak dari $A$ ke bidang itu
        (\cref{prop:g12:space:distance}).
  \item Apakah \omterm{def:g11:vect:vector}{vektor} $(1,-1,0)$ dan $(1,1,-2)$ \omterm{def:g11:scal:orthogonal}{ortogonal}?
  \item Bagaimana kedudukan bidang $x + y + z = 1$ dan
        $2x + 2y + 2z = 5$? Lalu kedudukan garis lewat $A$ yang diarahkan oleh
        $(1,1,0)$ terhadap bidang yang pertama?
  \item Daftarkanlah \omterm{prop:g10:coordgeom:midpoint}{titik tengah} keenam rusuk
        $[BC], [CD], [DH], [HE], [EF], [FB]$ pada kubus itu, lalu
        periksalah bahwa keenamnya memenuhi \omterm{def:g10:algebra:equation}{persamaan}
        $x + y + z = \frac32$.
\end{enumerate}

\medskip
\noindent\textbf{Bagian II --- \omterm{def:g10:numbers:interunion}{Irisan} segi enamnya.}
Misalkan $P$ bidang $x + y + z = \frac32$.
\begin{enumerate}[resume]
  \item Tunjukkan bahwa $P$ tegak lurus terhadap diagonal besar
        $(AG)$ dan lewat pusat kubusnya.
  \item Hitunglah jarak antara \omterm{prop:g10:coordgeom:midpoint}{titik tengah} yang berurutan pada
        pertanyaan 5 sepanjang \omterm{def:g10:numbers:interunion}{irisannya}: apa yang kamu peroleh?
  \item Hitunglah sudut segi banyak \omterm{def:g10:numbers:interunion}{irisannya} pada satu
        titik sudut (dua \omterm{def:g11:vect:vector}{vektor} rusuk yang berurutan dan satu \omterm{def:g12:space:dot}{hasil
        kali skalar}). Simpulkan: \omterm{def:g10:numbers:interunion}{irisannya} adalah \emph{segi enam
        beraturan}.
  \item Hitunglah keliling dan luas segi enamnya, lalu bandingkan
        luasnya dengan luas satu sisi kubusnya. (Sebagai
        catatan: \omterm{def:g10:numbers:interunion}{irisan} bidang terbesar pada kubus satuan adalah
        persegi panjang diagonal $1 \times \sqrt2$, yang berluas
        $\sqrt2$ --- jadi segi enam beraturan itu meraih perak.)
  \item Di manakah diagonal $(AG)$ menembus $P$? Periksalah bahwa
        titik tembusnya adalah \omterm{prop:g10:coordgeom:midpoint}{titik tengah} $[AG]$.
  \item Geserlah \omterm{def:g10:numbers:interunion}{irisannya}: lukiskan \omterm{def:g10:numbers:interunion}{irisan}
        $x + y + z = c$ ketika $c$ tumbuh dari $0$ sampai $3$ ---
        hitunglah \omterm{def:g10:numbers:interunion}{irisannya} untuk $c = \frac12$ (segi banyak mana,
        seberapa besar?), lalu ceritakan perubahan bentuknya yang melewati
        segi enammu di $c = \frac32$.
  \item Mengapa tak ada bidang yang dapat memotong kubus menjadi segi banyak
        bersisi \emph{tujuh}? (Di manakah tiap sisi \omterm{def:g10:numbers:interunion}{irisannya}
        harus terletak?)
\end{enumerate}

\medskip
\noindent\textbf{Bagian III --- Bidang empat yang tersembunyi.}
\begin{enumerate}[resume]
  \item Hitunglah volume bidang empat sudut $ABDE$
        (alasnya, tingginya, dan rumus sepertiga yang diterima saja
        oleh jilid sekolah menengah pertama --- \cref{exo:g12:space:9}
        menghitung volume semacam itu secara umum).
  \item Keempat titik sudut $B$, $D$, $E$, $G$: hitunglah keenam
        jarak berpasangannya lalu simpulkan bahwa $BDEG$ adalah
        bidang empat yang \emph{beraturan}. Lalu perolehlah volumenya dengan
        mengurangkan bidang empat sudutnya dari kubus itu.
  \item Hitunglah jarak dari pusat kubusnya ke
        bidang $(BDE)$.
  \item Sudut intannya: hitunglah sudut antara
        diagonal $\vect{AG}$ dan $\vect{BH}$ pada kubus itu. Pelurusnya,
        yang kira-kira $109.5^\circ$, adalah sudut antara
        ikatan pada intan dan metana --- mengapa empat
        pasangan elektron di sekeliling atom karbon memilih sudut
        bidang empat pada pertanyaan 14?
\end{enumerate}

\medskip
\noindent\textbf{Bagian IV --- Sungguh-sungguh tiga dimensi.}
\begin{enumerate}[resume]
  \item Carilah tempat garis lewat $B(1,0,0)$ yang diarahkan oleh
        $(0,1,1)$ menembus bidang segi enamnya, $P$.
  \item Tunjukkan bahwa garis $(AB)$ dan $(EG)$ tidak bertemu dan tidak pula
        sejajar: itulah garis \emph{bersilangan}, gejala yang
        tak mungkin terjadi pada sebuah bidang.
  \item Hitunglah proyeksi \omterm{def:g11:scal:orthogonal}{ortogonal} $B$ pada $P$
        (berjalanlah dari $B$ sepanjang normalnya sampai ke bidangnya), lalu
        periksalah bahwa titikmu memenuhi \omterm{def:g10:algebra:equation}{persamaan} $P$.
  \item Penutup --- perlengkapan sang juru ukur ruang: garis parametrik
        untuk menembus, \omterm{def:g12:space:normal}{vektor normal} untuk sudut dan jarak,
        \omterm{def:g10:algebra:equation}{persamaan} bidang untuk \omterm{def:g10:numbers:interunion}{irisan}; dan kubus sebagai
        tempat bermain yang di sana semuanya menghasilkan segi enam, sebuah
        bidang empat beraturan, sudut intan dan sepasang garis
        bersilangan. Masing-masing satu kalimat.

\end{enumerate}
\end{problem}
